\scanpage{95}
savoir les ``invariants birationnels fondamentaux'' \deleted{\unsure{de}} \inserted{relativement à $\mathbb F_p$,} envisagés dans les notes sur les \uline{Motifs}.

\noindent\makebox[0pt][r]{\rule[-3\baselineskip]{.35pt}{3.9\baselineskip}\hspace{2pt}\rule[-3\baselineskip]{.35pt}{3.9\baselineskip}\hspace{.55em}}Il semblerait opportun d'étudier \deleted{\unsure{si}} \inserted{le cas de coefficients $F$ généraux et de voir si,} pour un morphisme $f:X\longrightarrow Y$ disons lisse \deleted{\unsure{projectif}}propre, \inserted{$F$ \unsure{cst|ct} tordu sur $X$,} la connaissance de \deleted{\unsure{$H^i$|$H$}} $L(X,F)$ modulo $L(f^{-1}(Y'),F)$, pour $Y'$ \unsure{var.|var|variant} dans $Y$, permet de récupérer les invariants birationnels fondamentaux de $(X_K,F\mid X_K)/K$ [$K$ corps des fonctions de $Y$]. \uline{Dans quelle mesure faut-il faire les hypothèses $f$ lisse propre}, et $F$ \uline{ct tordu} pour obtenir des énoncés de cette nature ??
% archive sheet number: 94
